\documentclass[12pt,a4paper]{report}
\usepackage[utf8]{inputenc}
\usepackage[T1]{fontenc}
\usepackage[french]{babel}
\usepackage{geometry}
\usepackage{array}
\usepackage{longtable}
\usepackage{xcolor}
\usepackage{booktabs}
\usepackage{makecell}
\geometry{margin=2cm}

\definecolor{headerblue}{RGB}{30, 80, 160}

% Custom column type for left-aligned paragraph
\newcolumntype{L}[1]{>{\raggedright\arraybackslash}p{#1}}

\begin{document}

% ═══════════════════════════════════════════════════════════════════════════════
%  TABLE 1 — CAS D'UTILISATION : S'INSCRIRE
% ═══════════════════════════════════════════════════════════════════════════════
\begin{center}
\renewcommand{\arraystretch}{1.5}
\begin{longtable}{|L{4cm}|L{11.5cm}|}
\hline
\rowcolor{headerblue}
\multicolumn{2}{|l|}{\color{white}\textbf{Table -- Description textuelle du cas d'utilisation \guillemotleft~S'inscrire~\guillemotright}} \\
\hline
\textbf{ID}               & UC-1 \\
\hline
\textbf{Nom}              & S'inscrire \\
\hline
\textbf{Acteurs primaires} & Utilisateur (non enregistré) \\
\hline
\textbf{Description}      & L'utilisateur souhaite créer un compte SoulLink afin d'accéder à la plateforme. Il fournit ses informations personnelles, choisit un identifiant unique (handle) et reçoit un code OTP pour vérifier son e-mail ou son numéro de téléphone. \\
\hline
\textbf{Préconditions}    &
\begin{itemize}
  \item L'utilisateur ne possède pas encore de compte sur la plateforme.
  \item L'utilisateur doit avoir au moins 18 ans (validé automatiquement par le système).
  \item L'utilisateur doit fournir au moins un e-mail \textbf{ou} un numéro de téléphone.
\end{itemize} \\
\hline
\textbf{Scénario nominal} &
\begin{enumerate}
  \item L'utilisateur accède à la page d'inscription.
  \item Le système affiche le formulaire d'inscription (nom affiché, handle, e-mail ou téléphone, mot de passe, date de naissance, ville, pays).
  \item L'utilisateur remplit les champs obligatoires et valide le formulaire.
  \item Le système vérifie que les données saisies respectent les contraintes de validation (mot de passe $\geq$ 8 caractères, handle alphanumérique, âge $\geq$ 18 ans).
  \item Le système vérifie que l'e-mail, le téléphone et le handle ne sont pas déjà utilisés.
  \item Le système crée le compte avec le statut \texttt{PENDING\_VERIFICATION}.
  \item Le système génère un code OTP à 6 chiffres (valide 15 minutes) et l'envoie par e-mail ou par SMS.
  \item Le système génère un \textit{access token} et un \textit{refresh token} et connecte automatiquement l'utilisateur.
  \item Le système redirige l'utilisateur vers l'écran de vérification OTP.
\end{enumerate} \\
\hline
\textbf{Scénarios alternatifs} &
\textbf{4.1 : Données manquantes ou invalides} \\
& \quad 4.1.1 Le système détecte qu'un champ obligatoire est absent ou invalide. \\
& \quad 4.1.2 Le système affiche un message d'erreur précis (ex. : « Le mot de passe doit comporter au moins 8 caractères »). \\
& \quad 4.1.3 Retour à l'étape 2 du scénario nominal. \\[4pt]
& \textbf{4.2 : L'utilisateur n'a pas encore 18 ans} \\
& \quad 4.2.1 Le système détecte que la date de naissance indique un âge inférieur à 18 ans. \\
& \quad 4.2.2 Le système affiche le message : « Vous devez avoir au moins 18 ans pour rejoindre SoulLink ». \\
& \quad 4.2.3 Retour à l'étape 2 du scénario nominal. \\[4pt]
& \textbf{5.1 : E-mail, téléphone ou handle déjà utilisé} \\
& \quad 5.1.1 Le système détecte qu'un compte existant utilise déjà cet e-mail, ce téléphone ou ce handle. \\
& \quad 5.1.2 Le système affiche le message : « Un utilisateur avec cet e-mail, ce téléphone ou ce handle existe déjà ». \\
& \quad 5.1.3 Retour à l'étape 2 du scénario nominal. \\[4pt]
& \textbf{5.2 : Identité faciale déjà enregistrée (si descripteur facial fourni)} \\
& \quad 5.2.1 Le système compare le descripteur facial soumis avec ceux déjà stockés en base. \\
& \quad 5.2.2 Si la distance euclidienne est inférieure à 0,45, le système refuse l'inscription. \\
& \quad 5.2.3 Le système affiche : « Cette identité faciale est déjà enregistrée sur un autre compte ». \\
& \quad 5.2.4 Retour à l'étape 2 du scénario nominal. \\[4pt]
& \textbf{7.1 : Échec d'envoi du code OTP} \\
& \quad 7.1.1 Le service d'e-mail ou de SMS rencontre une erreur technique. \\
& \quad 7.1.2 Le compte est néanmoins créé en base de données. \\
& \quad 7.1.3 Le système invite l'utilisateur à redemander l'envoi du code depuis son profil. \\
\hline
\textbf{Postconditions} &
\begin{itemize}
  \item Le compte est créé avec le statut \texttt{PENDING\_VERIFICATION}.
  \item Un code OTP est enregistré en base (valide 15 minutes) et envoyé à l'utilisateur.
  \item L'utilisateur est connecté automatiquement (tokens générés) en attente de vérification.
\end{itemize} \\
\hline
\end{longtable}
\end{center}

\vspace{1cm}

% ═══════════════════════════════════════════════════════════════════════════════
%  TABLE 2 — CAS D'UTILISATION : SE CONNECTER
% ═══════════════════════════════════════════════════════════════════════════════
\begin{center}
\renewcommand{\arraystretch}{1.5}
\begin{longtable}{|L{4cm}|L{11.5cm}|}
\hline
\rowcolor{headerblue}
\multicolumn{2}{|l|}{\color{white}\textbf{Table -- Description textuelle du cas d'utilisation \guillemotleft~Se connecter~\guillemotright}} \\
\hline
\textbf{ID}               & UC-2 \\
\hline
\textbf{Nom}              & Se connecter \\
\hline
\textbf{Acteurs primaires} & Utilisateur (Administrateur, Modérateur, Utilisateur standard) \\
\hline
\textbf{Description}      & L'utilisateur souhaite accéder à son espace personnel SoulLink. Il saisit son identifiant (e-mail ou téléphone) et son mot de passe, ou utilise la reconnaissance faciale. Le système valide les informations et délivre un jeton d'accès. \\
\hline
\textbf{Préconditions}    &
\begin{itemize}
  \item L'utilisateur possède un compte valide sur la plateforme.
  \item Le compte n'est pas \texttt{BANNED} ni \texttt{SUSPENDED}.
\end{itemize} \\
\hline
\textbf{Scénario nominal \newline (connexion par identifiant)} &
\begin{enumerate}
  \item L'utilisateur accède à la page de connexion.
  \item Le système affiche le formulaire de connexion (identifiant e-mail/téléphone et mot de passe).
  \item L'utilisateur saisit son identifiant et son mot de passe, puis valide.
  \item Le système recherche le compte correspondant à l'identifiant fourni.
  \item Le système compare le mot de passe saisi avec le hash stocké en base (bcrypt).
  \item Le système génère un \textit{access token} (JWT) et un \textit{refresh token} (cookie HTTP-only, 7 jours).
  \item Le système enregistre la connexion dans l'historique (\textit{IP}, \textit{User-Agent}, horodatage).
  \item Le système met à jour la date de dernière connexion et le fuseau horaire de l'utilisateur.
  \item Le système redirige l'utilisateur vers son tableau de bord personnel.
\end{enumerate} \\
\hline
\textbf{Scénario nominal \newline (connexion faciale)} &
\begin{enumerate}
  \item L'utilisateur choisit l'option « Se connecter par reconnaissance faciale ».
  \item Le système active la caméra et capture le descripteur facial (vecteur à 128 dimensions).
  \item Le système compare le descripteur capturé avec l'ensemble des descripteurs stockés en base.
  \item Le système calcule la distance euclidienne entre le vecteur capturé et chaque vecteur stocké.
  \item Si la distance minimale est inférieure à 0,6, le système identifie l'utilisateur correspondant.
  \item Le système génère les jetons d'accès, met à jour \texttt{faceVerified} et \texttt{lastLoginAt}.
  \item Le système connecte l'utilisateur et le redirige vers son tableau de bord.
\end{enumerate} \\
\hline
\textbf{Scénarios alternatifs} &
\textbf{4.1 : Identifiant introuvable} \\
& \quad 4.1.1 Le système ne trouve aucun compte correspondant à l'e-mail ou au téléphone saisi. \\
& \quad 4.1.2 Le système affiche : « Identifiant ou mot de passe invalide ». \\
& \quad 4.1.3 Retour à l'étape 2 du scénario nominal. \\[4pt]
& \textbf{5.1 : Mot de passe incorrect} \\
& \quad 5.1.1 La comparaison bcrypt échoue. \\
& \quad 5.1.2 Le système affiche : « Identifiant ou mot de passe invalide ». \\
& \quad 5.1.3 Retour à l'étape 2 du scénario nominal. \\[4pt]
& \textbf{5.2 : Compte suspendu ou banni} \\
& \quad 5.2.1 Le système détecte que le statut du compte est \texttt{SUSPENDED} ou \texttt{BANNED}. \\
& \quad 5.2.2 Le système affiche : « Votre compte est suspendu / banni ». \\
& \quad 5.2.3 La connexion est refusée. \\[4pt]
& \textbf{5.3 : Tentatives excessives (rate limiting)} \\
& \quad 5.3.1 Le \textit{rate limiter} détecte plus de requêtes autorisées depuis la même IP. \\
& \quad 5.3.2 Le système bloque temporairement l'accès et affiche un message d'attente. \\
& \quad 5.3.3 L'utilisateur peut réessayer après la fenêtre de blocage. \\[4pt]
& \textbf{Facial — 5.4 : Visage non reconnu} \\
& \quad 5.4.1 La distance euclidienne minimale calculée est supérieure ou égale à 0,6. \\
& \quad 5.4.2 Le système affiche : « Visage non reconnu ». \\
& \quad 5.4.3 Le système propose à l'utilisateur de se connecter par identifiant/mot de passe. \\[4pt]
& \textbf{Facial — 5.5 : Aucune donnée faciale en base} \\
& \quad 5.5.1 Aucun utilisateur n'a encore enregistré de descripteur facial. \\
& \quad 5.5.2 Le système affiche : « Aucun utilisateur trouvé avec des données faciales ». \\
\hline
\textbf{Postconditions} &
\begin{itemize}
  \item L'utilisateur est authentifié et possède un \textit{access token} valide.
  \item Un \textit{refresh token} est stocké dans un cookie sécurisé (\texttt{HttpOnly}, \texttt{SameSite=strict}).
  \item La connexion est enregistrée dans l'historique (\texttt{LoginHistory}).
  \item La date de dernière connexion (\texttt{lastLoginAt}) est mise à jour en base.
\end{itemize} \\
\hline
\end{longtable}
\end{center}

\vspace{1cm}

% ═══════════════════════════════════════════════════════════════════════════════
%  TABLE 3 — CAS D'UTILISATION : RÉINITIALISER LE MOT DE PASSE
% ═══════════════════════════════════════════════════════════════════════════════
\begin{center}
\renewcommand{\arraystretch}{1.5}
\begin{longtable}{|L{4cm}|L{11.5cm}|}
\hline
\rowcolor{headerblue}
\multicolumn{2}{|l|}{\color{white}\textbf{Table -- Description textuelle du cas d'utilisation \guillemotleft~Réinitialiser le mot de passe~\guillemotright}} \\
\hline
\textbf{ID}               & UC-3 \\
\hline
\textbf{Nom}              & Réinitialiser le mot de passe \\
\hline
\textbf{Acteurs primaires} & Utilisateur (Administrateur, Modérateur, Utilisateur standard) \\
\hline
\textbf{Description}      & L'utilisateur ayant oublié son mot de passe souhaite en définir un nouveau. Il fournit son identifiant (e-mail ou téléphone), reçoit un code OTP de réinitialisation valide 15 minutes, puis saisit et confirme un nouveau mot de passe. \\
\hline
\textbf{Préconditions}    &
\begin{itemize}
  \item L'utilisateur possède un compte actif enregistré avec un e-mail ou un numéro de téléphone valide.
  \item L'utilisateur a accès à son e-mail ou à son téléphone pour recevoir le code OTP.
\end{itemize} \\
\hline
\textbf{Scénario nominal} &
\begin{enumerate}
  \item L'utilisateur clique sur « Mot de passe oublié » sur la page de connexion.
  \item Le système affiche un champ de saisie pour l'identifiant (e-mail ou numéro de téléphone).
  \item L'utilisateur saisit son identifiant et soumet la demande.
  \item Le système recherche le compte correspondant à l'identifiant fourni.
  \item Le système génère un code OTP à 6 chiffres et crée une entrée de vérification de type \texttt{PASSWORD\_RESET} valide 15 minutes.
  \item Le système envoie le code OTP par e-mail (si identifiant = e-mail) ou par SMS (si identifiant = téléphone).
  \item Le système affiche le message : « Si un compte existe, un code de réinitialisation a été envoyé ».
  \item L'utilisateur saisit le code OTP reçu, le nouvel mot de passe (minimum 8 caractères) et confirme.
  \item Le système vérifie que le code OTP correspond et n'est pas expiré.
  \item Le système hache le nouveau mot de passe (bcrypt) et met à jour le compte en base.
  \item Le système marque le code OTP comme utilisé (\texttt{verified = true}).
  \item Le système affiche le message de succès : « Mot de passe réinitialisé avec succès ».
  \item Le système redirige l'utilisateur vers la page de connexion.
\end{enumerate} \\
\hline
\textbf{Scénarios alternatifs} &
\textbf{4.1 : Identifiant introuvable} \\
& \quad 4.1.1 Le système ne trouve aucun compte associé à l'identifiant fourni. \\
& \quad 4.1.2 Le système affiche tout de même : « Si un compte existe, un code a été envoyé » (pour des raisons de sécurité, aucune information sur l'existence du compte n'est divulguée). \\
& \quad 4.1.3 Aucun code n'est généré ni envoyé. \\[4pt]
& \textbf{4.2 : Rate limiting activé (tentatives excessives)} \\
& \quad 4.2.1 Le \textit{rate limiter} détecte un nombre excessif de requêtes depuis la même IP. \\
& \quad 4.2.2 Le système bloque temporairement les nouvelles demandes et affiche un message d'attente. \\[4pt]
& \textbf{9.1 : Code OTP incorrect} \\
& \quad 9.1.1 Le code saisi par l'utilisateur ne correspond à aucun enregistrement valide de type \texttt{PASSWORD\_RESET}. \\
& \quad 9.1.2 Le système affiche : « Code de réinitialisation invalide ou expiré ». \\
& \quad 9.1.3 Retour à l'étape 8 du scénario nominal. \\[4pt]
& \textbf{9.2 : Code OTP expiré} \\
& \quad 9.2.1 La date d'expiration du code enregistré (\texttt{expiresAt}) est antérieure à l'heure actuelle. \\
& \quad 9.2.2 Le système affiche : « Le code de réinitialisation a expiré ». \\
& \quad 9.2.3 Le système invite l'utilisateur à recommencer la procédure depuis l'étape 1. \\[4pt]
& \textbf{9.3 : Nouveau mot de passe trop court} \\
& \quad 9.3.1 Le système détecte que le nouveau mot de passe comporte moins de 8 caractères. \\
& \quad 9.3.2 Le système affiche : « Le mot de passe doit comporter au moins 8 caractères ». \\
& \quad 9.3.3 Retour à l'étape 8 du scénario nominal. \\
\hline
\textbf{Postconditions} &
\begin{itemize}
  \item Le mot de passe de l'utilisateur est mis à jour en base de données (hashé avec bcrypt).
  \item Le code OTP de type \texttt{PASSWORD\_RESET} est marqué comme utilisé (\texttt{verified = true}).
  \item L'utilisateur peut désormais se connecter avec son nouveau mot de passe.
\end{itemize} \\
\hline
\end{longtable}
\end{center}

\end{document}
